\subsection{甲状腺功能异常与性欲：被忽视的"代谢--情绪--性"三角}

\noindent\textit{【编者注】}甲状腺是人体的"代谢司令部"，其激素几乎影响每一个器官——包括性腺与大脑的性欲中枢。甲状腺功能异常在普通人群患病率不低（女性显著高于男性），却常在"性欲低下查因"时被漏问。本节补齐这一环。

\vspace{0.5em}
\noindent\textbf{甲状腺功能亢进（甲亢）对性的影响}：

\begin{itemize}
	\item \textbf{早期可表现为性欲"反常升高"}：代谢加快、情绪亢奋，部分患者性欲增强——但这不是"身体好"，而是疾病信号，常伴随心慌、手抖、消瘦、易怒与失眠
	\item \textbf{随病程进展多转为性功能减退}：慢性消耗与激素紊乱之下，男性可出现 ED、乳房发育（男性乳腺增生），女性常见月经稀发或紊乱、性唤起下降
	\item \textbf{情绪与关系层面}：易激惹、焦虑会让伴侣间的亲密氛围变得紧张，"脾气变大"往往先于"性欲变化"被注意到
\end{itemize}

\vspace{0.5em}
\noindent\textbf{甲状腺功能减退（甲减）对性的影响}：

\begin{itemize}
	\item \textbf{性欲低下是最常见的表现}：代谢整体"降速"，伴随乏力、怕冷、嗜睡、体重增加与情绪低落——性欲被"连根抑制"
	\item \textbf{激素层面的连锁反应}：甲减可引起高泌乳素血症，进一步抑制下丘脑--垂体--性腺轴，导致男性 ED 与精子质量下降、女性月经紊乱与排卵障碍、阴道干涩与性交不适
	\item \textbf{容易被误判}：甲减的症状（疲劳、性欲低、情绪低）与抑郁高度重叠，不少人被当作"单纯的抑郁"或"累了"而拖延多年
\end{itemize}

\vspace{0.5em}
\noindent\textbf{识别线索与就医路径}：

\begin{enumerate}
	\item 出现不明原因的性欲显著变化，尤其伴随心慌/怕冷、体重骤变、月经紊乱、脱发、颈部增粗或情绪异常时，建议查一次\textbf{甲状腺功能五项}（TSH、FT3、FT4 及抗体），花费不高、抽一次血即可
	\item 确诊后以\textbf{治疗原发病为主}：甲亢经规范治疗（药物、放射性碘或手术）后性功能多随代谢正常化而恢复；甲减患者规律补充左甲状腺素、把 TSH 控制在目标范围后，性欲低下多数可逆
	\item 性功能问题持续存在者，再由内分泌科与男科/妇科协同排查其他原因，避免"一把抓"
\end{enumerate}

\begin{tcolorbox}[colback=pink!5!white,colframe=red!50!black,title=贴心小叮咛]
甲状腺问题引起的性功能变化\textbf{绝大多数是可逆的}。它最常被忽略的原因，是大家不会把"床上的事"和脖子上的蝴蝶形小腺体联系起来。下次体检时多问一句"要不要查个甲状腺"，可能就解开了一个困扰很久的谜。
\end{tcolorbox}

